\chapter{KESIMPULAN DAN SARAN}

\section{Kesimpulan}
Berdasarkan hasil perancangan, pengujian, dan pembahasan yang telah diuraikan pada bab-bab sebelumnya, kesimpulan dari penelitian ini adalah sebagai berikut:
\begin{enumerate}[label=\arabic*.]
    \item Sistem tertanam berbasis Raspberry Pi 5 berhasil dirancang dan diimplementasikan sebagai pusat komputasi tepi (\textit{edge computer}) yang mengintegrasikan kamera web Logitech, \textit{framework} MediaPipe, Python 3.9, serta dua robot manipulator 4-DOF (meArm dan Dobot Magician) secara lokal tanpa bergantung pada komputasi awan. \textit{Pipeline} visi komputer berjalan stabil pada perangkat Raspberry Pi 5 untuk kedua mode operasi.
    \item Mode pengenalan Sistem Isyarat Bahasa Indonesia (SIBI) pada meArm berhasil mengklasifikasikan 26 abjad dengan akurasi keseluruhan sebesar \textbf{96,73\%} (503 prediksi benar dari 520 percobaan). Model \textit{Deep Neural Networks} (DNN) berbasis 63 fitur koordinat spasial \textit{landmark} tangan mampu mengenali 22 dari 26 kelas abjad dengan akurasi sempurna (100,0\%), dengan galat klasifikasi terbatas pada huruf J, M, N, dan Y akibat kemiripan topologi gestur.
    \item Mode teleoperasi pada Dobot Magician berhasil memetakan gerak tangan operator ke ruang kerja robot dengan kesesuaian 6 dari 6 arah pemetaan (valid), galat spasial rata-rata sebesar \textbf{1,90 mm}, dan waktu respons rata-rata sebesar \textbf{87,25 ms}. Penerapan filter \textit{exponential smoothing} ($\alpha = 0{,}2$) dan ambang batas gerak 5,0 mm berhasil mengeliminasi getaran mikro (\textit{jitter}) pada motor penggerak robot.
\end{enumerate}

\section{Saran}
Saran untuk pengembangan penelitian selanjutnya adalah sebagai berikut:
\begin{enumerate}[label=\arabic*.]
    \item Menambah jumlah kelas dan variasi data SIBI dari beberapa pengguna agar sistem lebih tahan terhadap variasi bentuk tangan.
    \item Menggunakan metode klasifikasi yang lebih kuat, seperti \textit{support vector machine}, \textit{random forest}, atau jaringan saraf ringan yang tetap dapat dijalankan pada Raspberry Pi 5.
    \item Menambahkan filter temporal pada \textit{landmark} tangan agar gerak robot lebih halus dan tidak bergetar.
    \item Melakukan kalibrasi ruang kerja yang lebih presisi antara kamera dan Dobot Magician.
    \item Menambahkan antarmuka pengguna untuk memilih mode, memantau hasil klasifikasi, dan menyimpan data pengujian secara otomatis.
\end{enumerate}
